\chapter{Preliminaries and Guide to the Notes}\label{chap:preliminaries}

\section{Purpose and point of view}

Graduate algebra is not chiefly a collection of classifications.  It is a
language for recognizing when apparently different problems have the same
form. A quotient group, a quotient ring, and a quotient module all arise by
collapsing a chosen subobject to the trivial one: a normal subgroup becomes
the identity in a group, while an ideal or submodule becomes zero. An action describes symmetry; a
polynomial in an operator turns linear algebra into module theory; and a
Galois group records the symmetries of a field extension.  These notes develop
that common language with an eye toward the places it will be used next.

The intended reader has strong high-school or early undergraduate mathematical
maturity: comfort with sets, functions, elementary proof, and familiar number
systems such as \(\mathbb Z\), \(\mathbb Q\), and \(\mathbb R\).  No prior
course in abstract algebra is assumed.  Definitions begin from that elementary
entry point, examples arrive before abstraction, and proofs introduce each new
move before relying on it.  The destination is nevertheless graduate algebra:
the resulting core is designed for later commutative algebra, homological
algebra, representation theory, algebraic geometry, and algebraic number
theory.  The last chapters give a complete first treatment of classical Galois
theory.

This cumulative policy is deliberate.  Early chapters do not appeal to an
unseen ``standard fact from algebra''; later chapters may use a result only
after it has been established here.  Readers who have already studied a first
course can move more quickly through the examples, but they are not the
prerequisite for understanding the logical development.

Readers who would like additional preparation in mathematical language,
proof-writing, foundational constructions, and the number systems may consult
Chapter~2 of \emph{Lecture Notes on Mathematical Analysis} \cite{MA}.
Its discussions of sets, functions, proof conventions, induction, and the
construction and basic properties of \(\mathbb N,\mathbb Z,\mathbb Q,\mathbb R\)
can help bridge to the proof-based style used here.  This is useful
preparation, not a prerequisite: all algebraic material required in these
notes is developed within the manuscript.

Two habits are central.  First, definitions are treated as constructions with
universal features, rather than as names for collections of examples.  Second,
proofs are organized around reusable moves: pass to a quotient, use a universal
property, compare kernels and images, act on a set, reduce a statement to a
cyclic object, or use induction on a degree or a composition length.  The same
moves will recur throughout the book.

\section{The mathematical route}

Part~I establishes the group-theoretic language of normal subgroups, quotients,
actions, and composition factors.  Sylow theory is included for its structural
consequences, while lengthy classifications of groups of individual small
orders are not.  Part~II develops division rings, ideals, quotients, group
rings, and localization, then the chain
of implications
\[
  \text{Euclidean domain}\quad\Longrightarrow\quad\text{PID}
  \quad\Longrightarrow\quad\text{UFD}.
\]
It also develops Noetherian finiteness, Hilbert's basis theorem, formal power
series, one- and several-variable polynomial rings, finite-field applications,
a short quadratic-reciprocity perspective, and an introduction to Gr\"obner
bases.  Localization is included because it is the first algebraic operation
that makes a global ring look local.

Part~III develops four connected chapters. Chapter~5 introduces modules
together with the finite-dimensional linear algebra needed later: bases,
dimension, matrices, determinants, eigenvalues, diagonalization, duality,
and a modest introduction to bilinear and quadratic forms.
Chapter~6 proves a strengthening theorem for submodules of finite free
modules over a PID, producing compatible ambient coordinates one
invariant-factor direction at a time. It then derives the PID structure
theorem. Uniqueness in elementary-divisor form is recovered from intrinsic
prime-primary layers, while uniqueness in invariant-factor form is also
proved directly by annihilator, cyclic cancellation, and induction. Smith
normal form gives the computational
matrix manifestation before the chapter regards an operator as an
\(F[x]\)-module, leading to minimal polynomials and rational and Jordan
canonical forms. Chapter~7 adds inner-product geometry, adjoints, spectral
theory, and singular value decomposition. Chapter~8 introduces tensor
products, exterior powers, and exterior algebra: the module constructions
that linearize bilinear and alternating multilinear maps and prepare the way
for later commutative and homological algebra. Thus vector spaces are
modules over fields, abelian groups are modules over \(\mathbb Z\), and a
linear operator on an \(F\)-vector space makes that space into an
\(F[x]\)-module. The structure theorem for finitely generated modules over a
PID therefore simultaneously explains finitely generated abelian groups and
the rational canonical form of a linear map. Jordan form appears as the
split-polynomial refinement of the same idea.

Part~IV studies field extensions, their embeddings, separability, normality,
and the correspondence between intermediate fields and subgroups of a Galois
group.  Finite fields and cyclotomic fields supply recurring examples, while
direct and inverse limits explain how finite Galois symmetries assemble into
profinite and absolute Galois groups.  Chapter~11 includes a brief local
bridge through point-set topology and topological groups, solely to make the
profinite and infinite-Galois language self-contained; Tychonoff's theorem is
used there as an external input.  It closes with a non-prerequisite survey of
the later directions opened by all four parts. The later discussion of
radicals illustrates how questions about polynomial formulas are controlled by
group-theoretic properties.

\section{How to read a proof}

Every result used later is stated with its hypotheses and proved before its
first essential use.  A proof normally begins by identifying its idea: for
example, the first isomorphism theorem compares an object with the quotient by
the kernel of a map, whereas the orbit--stabilizer theorem compares a group
action with a quotient by a stabilizer.  Readers should pause at these ideas,
not merely at the final calculation.

Some results are deliberately stated in a general form.  This is done only
when the extra generality clarifies later work.  For instance, the module
isomorphism theorems subsume their group and ring analogues, but the earlier
versions are still proved because quotient constructions must be understood in
their native settings.  Conversely, the classification of finitely generated
abelian groups is presented as a corollary of the PID module theorem rather
than as an unrelated classification.

\section{Exercises and completeness}

Exercises have two roles: they give essential practice, and they sometimes
contain theorems that deserve a place in the main narrative.  We do not follow
the source's typography as a measure of importance.  When an exercise supplies
a standard lemma, a useful converse, a structural construction, or a result
needed later, the result is promoted and proved in the text.  The companion
coverage audit records the exact source exercise, its location here, and the
reason for promotion.

At the end of each chapter, exercises are grouped as \emph{Basic},
\emph{Core}, \emph{Further}, and \emph{Looking Ahead}.  The first group tests
definitions; the second develops results expected of every reader; the third
extends the chapter; and the last points toward later subjects without being a
logical prerequisite.  A routine exercise remains an exercise even when its
solution is long.

\section{Conventions and notation}

Rings in these notes are associative and have a multiplicative identity; the
zero ring is allowed.  Ring homomorphisms preserve identity unless explicitly
stated otherwise, and a subring contains the same identity as its ambient
ring.  Terminology concerning rings is not completely uniform: some authors
allow rings without a multiplicative identity, and some do not require a
subring to contain the ambient identity.  The conventions used here will
always be stated explicitly when a nonunital object is needed.  Unless
explicitly qualified otherwise, an ideal means a two-sided ideal.  When
one-sided ideals are relevant, we explicitly say left ideal or right ideal;
in a commutative ring these notions coincide.
Modules are unital left modules unless the contrary is specified. For a general
group \(G\), the identity is \(e\), and for an additive group it is \(0\).
We retain \(1\) for concrete multiplicative groups and \(I_n\) for matrix
groups; \(\operatorname{id}_X\) denotes an identity map. The
notation \(H\leq G\) permits equality, while \(H\trianglelefteq G\) denotes a
normal subgroup.  If \(R\) is a ring and \(M\) an \(R\)-module, then
\(\operatorname{Ann}_R(M)\) is its annihilator.

For a regular \(n\)-gon, \(D_n\) denotes its full dihedral symmetry group, so
\(|D_n|=2n\).  Many authors write \(D_{2n}\) for this same group; the notation
in these notes records the number of polygon vertices rather than the group
order.  Thus \(D_4\) is the symmetry group of a square.

\begin{remark}[Looking Ahead]
The repeated pattern ``kernel, image, quotient'' is the beginning of the
language of exact sequences.  We will use exact sequences of modules in an
elementary form, but derived functors and category theory are not prerequisites
for these notes.
\end{remark}

\section{Elementary complex arithmetic}

The complex numbers extend the real numbers by an element \(i\) satisfying
\(i^2=-1\). Every complex number has a unique form
\[
z=a+bi,\qquad a,b\in\mathbb R.
\]
The numbers \(a=\operatorname{Re}(z)\) and \(b=\operatorname{Im}(z)\) are,
respectively, its real and imaginary parts. We add and multiply by collecting
real and imaginary parts and using \(i^2=-1\):
\[
(a+bi)+(c+di)=(a+c)+(b+d)i,
\quad
(a+bi)(c+di)=(ac-bd)+(ad+bc)i.
\]
For instance,
\[
(2+3i)+(1-4i)=3-i,\qquad
(2+3i)(1-4i)=14-5i.
\]

These definitions inherit the familiar algebraic laws of real-number
arithmetic. They are closed under addition and multiplication; addition and
multiplication are associative and commutative; \(0=0+0i\) is the additive
identity, \(1=1+0i\) is the multiplicative identity, and
\(-(a+bi)=-a-bi\) is the additive inverse. For example, coordinatewise
addition gives
\begin{align*}
((a+bi)+(c+di))+(p+qi)
&=(a+c+p)+(b+d+q)i\\
&=(a+bi)+((c+di)+(p+qi)).
\end{align*}
The same calculation gives commutativity and the additive identity and inverse
laws. Expanding both sides of \(z(w+u)=zw+zu\) and collecting real and
imaginary parts proves distributivity; associativity and commutativity of
multiplication follow similarly from the corresponding real-number laws.
The usual cancellation arguments prove uniqueness of the additive and
multiplicative identities and of additive inverses.

Conjugation is a useful device for converting a complex number into a real
number. Define
\[
\overline z=a-bi,\qquad |z|=\sqrt{a^2+b^2}.
\]
Thus for \(z=3+4i\), \(\overline z=3-4i\) and \(|z|=5\). Direct coordinate
calculations give
\[
\overline{z+w}=\overline z+\overline w,\qquad
\overline{zw}=\overline z\,\overline w,\qquad
\overline{\overline z}=z,
\]
and
\[
z\overline z=(a+bi)(a-bi)=a^2+b^2=|z|^2.
\]
Consequently
\[
|zw|^2=(zw)\overline{zw}=z\overline z\,w\overline w=|z|^2|w|^2,
\]
so \(|zw|=|z||w|\), since moduli are nonnegative.

If \(z=a+bi\ne0\), then \(a\) and \(b\) are not both zero, so
\(a^2+b^2>0\). Hence
\[
z^{-1}=\frac{a-bi}{a^2+b^2}=\frac{\overline z}{|z|^2}
\]
is defined, and \(zz^{-1}=z\overline z/|z|^2=1\). For example,
\[
(3+4i)^{-1}=\frac{3-4i}{25}.
\]
This inverse is unique: if \(u\) and \(v\) both invert \(z\), then
\[
u=u\cdot1=u(zv)=(uz)v=1\cdot v=v.
\]
These properties---together with the existence of multiplicative inverses for
every nonzero element---are exactly the familiar algebraic laws that will
later be summarized by saying that \(\mathbb C\) is a field.

Geometrically, \(a+bi\) is the point \((a,b)\) in the plane, and \(|z|\) is
its Euclidean distance from the origin. Thus the complex \(n\)-th roots of
unity
\[
\mu_n=\{z\in\mathbb C:z^n=1\}
=\left\{\cos\frac{2\pi k}{n}+i\sin\frac{2\pi k}{n}:0\leq k<n\right\}
\]
are the \(n\) equally spaced points on the unit circle. For example,
\[
\mu_2=\{1,-1\},\qquad \mu_4=\{1,i,-1,-i\}.
\]
They will later provide important group examples.

\section{Sources}

The primary mathematical reference is \cite{DF}.  The separate coverage audit
records editorial provenance and verified exercise references; it is not part
of the reader-facing source apparatus.
